\documentclass[10pt,a4paper]{amsart}
\usepackage[margin=19mm]{geometry}
\usepackage{amsmath,amssymb,mathtools}
\usepackage[scr=rsfs,cal=euler]{mathalfa}
\usepackage{xfrac}
\usepackage[unicode,hidelinks,bookmarks=false]{hyperref}
\newcommand{\ie}{i.e.\ }
\newcommand{\cf}{cf.\ }
\newcommand{\resp}{respectively\ }
\newcommand{\Subsection}{\subsection}
\providecommand{\nobreakdash}{\nobreak\hbox{-}}
\setlength{\emergencystretch}{2em}
\multlinegap0pt
\newtheorem{prop}{Proposition}
\newtheorem{lemma}{Lemma}
\newcommand{\calI}{\mathcal{I}}
\title{A new proof of the Artin--Springer theorem\\ in Schur index $2$}
\author{Anne Qu\'eguiner-Mathieu, Jean-Pierre Tignol}
\date{Source: arXiv:2504.16514v2 (CC BY 4.0)}
\begin{document}
\maketitle

\noindent\emph{Source and licence.} A new proof of the Artin--Springer theorem\\ in Schur index $2$. Version arXiv:2504.16514v2 (CC BY 4.0), distributed by arXiv under the Creative Commons Attribution 4.0 International licence, http://creativecommons.org/licenses/by/4.0/, as recorded in the arXiv licence metadata for this exact version and shown on the abstract page https://arxiv.org/abs/2504.16514v2. Full licence notice: \texttt{licenses/arxiv-2504.16514-CC-BY-4.0.txt}.\par\smallskip

\noindent\textit{Editorial scope.} Counted unit: the article's lemma lem:key (source0.tex lines 1510--1516). The article's complete original proof of it is delivered with it (source0.tex lines 1518--1615, one \texttt{proof} environment of 98 lines). No mathematics is written, repaired or re-derived: the statement and the proof are the article's own lines, in the article's own notation, and no retained line is substituted or re-flowed.  Retained uncounted as the source-local context the proof needs: lines 1220--1232; lines 1387--1410.  Nothing was omitted from any retained span, so the package carries no omission record.  No retained line contains a citation, so the wrapper carries no bibliography and no bibliography entry.  No retained line contains a figure environment, an image-inclusion command or drawing code, so no image is delivered and the package declares no asset dependency.  Editorial formatting only, in addition: an AMS \texttt{amsart} preamble that restates the article's own declarations and macros used by the retained lines, namely the environments \texttt{prop}, \texttt{lemma} on separate counters, together with the standard packages those lines need.  The retained environments print in this document as Proposition 1, Lemma 1 in order of appearance, whereas the article prints the same items with the numbering of its own full document.  The complete native source of the article as distributed in the arXiv e-print for this exact version is bundled unchanged as \texttt{sources/176954/source0.tex}.
\begin{prop}
  \label{prop:degI}
  For $d\in\mathbb{N}$, let $\calI_{\leq d} =
  \{\xi\in\calI\mid 
  \deg(\xi)\leq d\}$. Then
  \[
    \calI_{\leq0}=\{0\} \qquad\text{and}\qquad
    \calI_{\leq d} = \Bigl\{\sum_{i=0}^{d-1}x_it^i\varepsilon \mid
    x_i\in Q\Bigr\}
    =\Bigl\{\sum_{z=-d}^{d-1} y_zt^z\varepsilon \mid
    y_z\in K\Bigr\}\qquad\text{for $d\geq1$.}
  \]
\end{prop}

\begin{lemma}
  \label{lem:deg}
  For $d\in\mathbb{N}$,
  \[
    V\otimes_Q\calI_{\leq d+1} = \Bigl\{
    \sum_{i=0}^d v_i\otimes t^i\varepsilon\mid v_0,\,\ldots,\,v_d\in
    V\Bigr\}.
  \]
  Let $d$, $e\in\mathbb{N}$. If
  \[
    v=(v_{d}\otimes t^{d}\varepsilon) + v'\in V\otimes_Q\calI_{\leq d+1}
    \qquad\text{and}\qquad
    w = (w_{e}\otimes t^{e}\varepsilon) + w'\in V\otimes_Q\calI_{\leq e+1}
  \]
  with $v_{d}$, $w_{e}\in V$ and $v'\in V\otimes_Q\calI_{\leq
    d}$, $w'\in V\otimes_Q\calI_{\leq e}$, then
  \[
    \deg b_q(v,w) \leq d+e+1 \qquad\text{and}\qquad
    \deg q(v) \leq 2d+1.
  \]
  Moreover, if $\deg b_q(v,w)< d+e+1$, then $h_s(v_{d},
  w_{e})\in K$. If $\deg q(v) <2d+1$, then
  $s(v_{d},v_{d})\in K$.
\end{lemma}

\begin{lemma}
  \label{lem:key}
  Let $v\in V$ and $\xi\in\calI$ be such that $\deg(\xi)=d+1\geq2$. If
  $\deg q(v\otimes\xi)\leq2d-1$, then
  $q(v\otimes\xi) = q(v\otimes\eta)$ for some
  $\eta\in\calI_{\leq d}$.
\end{lemma}

\begin{proof}
  If $q(v\otimes\xi)=0$, the lemma holds with $\eta=0$. For the rest
  of the proof, we assume $q(v\otimes\xi)\neq0$.
  By Proposition~\ref{prop:degI} we may write
  $\xi=(x_{d}t^{d}+\cdots+x_0)\varepsilon$ with $x_{d}$,
  \ldots, $x_0\in Q$ and $x_{d}\neq 0$. Substituting $vx_d$ for $v$,
  we may (and will) assume $x_d=1$. Lemma~\ref{lem:deg} then yields
  $s(v,v)\in K$. If $s(v,v)\in k$, then
  \[
    q(v\otimes\xi) = a(\xi,s(v,v)\xi) = s(v,v) a(\xi,\xi)=0.
  \]
  This case is excluded since
  $q(v\otimes\xi)\neq0$, hence $s(v,v)\in K\setminus k$.

  By Proposition~\ref{prop:degI} we
  have
  \[
    (x_{d-2}t^{d-2}+\cdots+x_0)\varepsilon = (y_{d-2}t^{d-2}+\cdots +
    y_{-d+1}t^{-d+1})\varepsilon
  \]
  for some $y_{d-2}$, \ldots, $y_{-d+1}\in K$. We prove below:
  \medbreak

  \noindent\textit{Claim:} $x_{d-1}$ lies in $K$.
  \medbreak

  Assuming the claim, we complete the proof as follows: we write
  \[
    \xi=z\varepsilon \qquad\text{with}\qquad
    z=t^d+x_{d-1}t^{d-1}+y_{d-2}t^{d-2}+\cdots+y_{-d+1}t^{-d+1} \in
    K[t,t^{-1}].
  \]
  Since $s(v,v)\in K$, we get, by definition of $q$ and $a$,
  \begin{equation}
    \label{eq:key}
    q(v\otimes\xi) = a(\xi,s(v,v)\xi) = u\bigl(\iota(z)s(v,v)z -
    z\iota(s(v,v)z)\bigr) =
    uz\iota(z)\bigl(s(v,v)-\iota(s(v,v))\bigr).
  \end{equation}
  The rightmost expression does not change when we change $z$ into
  $\iota(z)$, 
  hence $q(v\otimes\xi) = q(v\otimes\iota(z)\varepsilon)$. 
  The lemma will follow because
  \[
    \iota(z)\varepsilon=(b^dt^{-d}+b^{d-1}\iota(x_{d-1})t^{1-d} +
    b^{d-2}\iota(y_{d-2}) t^{2-d} + \cdots +
    b^{-d+1}\iota(y_{-d+1})t^{d-1})\varepsilon \in \calI_{\leq d}.
  \]
  \medbreak

  To prove the claim,
  let $\xi_0=(x_{d-2}t^{d-2}+\cdots+x_0)\varepsilon\in \calI_{\leq
    d-1}$, so $\xi=(t^d+x_{d-1}t^{d-1})\varepsilon+\xi_0$ and
  \[
    q(v\otimes\xi) = q(v\otimes(t^d+x_{d-1}t^{d-1})\varepsilon) +
    b_q(v\otimes(t^d+x_{d-1}t^{d-1})\varepsilon, v\otimes\xi_0) +
    q(v\otimes\xi_0).
  \]
  Lemma~\ref{lem:deg} shows that the middle
  term on the right side has degree at most $2d-1$ and the rightmost
  term has degree at most $2d-3$. Therefore, the hypothesis that $\deg
  q(v\otimes\xi)\leq 2d-1$ implies
  \begin{equation}
    \label{eq:key2}
    \deg q(v\otimes(t^d+x_{d-1}t^{d-1})\varepsilon) \leq
    2d-1.
  \end{equation}
  Now, we have
  \[
    q(v\otimes(t^d+x_{d-1}t^{d-1})\varepsilon)=q(v\otimes
    t^d\varepsilon) + b_q(v\otimes t^d\varepsilon, v\otimes
    x_{d-1}t^{d-1}\varepsilon) + q(v\otimes x_{d-1}t^{d-1}\varepsilon)
  \]
  and Lemma~\ref{lem:deg} shows that the rightmost term has degree at
  most $2d-1$. Moreover, since $s(v,v)\in K\setminus
  k$, the same computation as for~\eqref{eq:key} yields
  \[
    q(v\otimes t^d\varepsilon) = ub^d\bigl(s(v,v)-\iota(s(v,v))\bigr)
    \in K^\times,
  \]
  hence $\deg q(v\otimes t^d\varepsilon) = 0$. Therefore,
  \eqref{eq:key2} implies
  \[
    \deg b_q(v\otimes t^d\varepsilon,v\otimes
    x_{d-1}t^{d-1}\varepsilon)\leq 2d-1,
    \qquad\text{i.e.,}\quad
    \deg a(t^d\varepsilon, h_s(v,vx_{d-1})t^{d-1}\varepsilon) \leq
    2d-1, 
  \]
  hence $h_s(v,\,vx_{d-1})\in K$ by
  Lemma~\ref{lem:deg}. But
  \[
    h_s(v,\,vx_{d-1}) = h_s(v,\,v)x_{d-1} = \bigl(s(v,v) -
    \overline{s(v,v)}\bigr) x_{d-1}
  \]
  and $s(v,v)-\overline{s(v,v)}\in K^\times$ because $s(v,v)\in
  K\setminus k$, hence $x_{d-1}\in K$. This completes the proof.
\end{proof}

\end{document}
